% problem-000143 2 含磷小分子反应
\section*{2. 含磷小分子反应}
\textit{下列为分问。}

\subsection*{2-1}
写出在−35 °C 时，将 $\mathrm { P H } _ { 3 }$ ⽓体通过⾦属钠的液氨溶液制备 $\mathrm { N a P H _ { 2 } }$ 的化学反应⽅程式及题中NaPH 与CO反应的化学反应⽅程式。

\subsection*{2-2}
画出化合物X的结构简式（溶剂分⼦⽤dme表⽰）。

\subsection*{2-3}
2011年Cummins等⼈报道了常温常压下，以⼄醚为溶剂，通过 $\mathrm { N a [ \{ B ( C _ { 6 } F _ { 5 } ) _ { 3 } \} P N b ( N [ N p ] A r ) }$ _{3}]与等摩尔数的 $\mathrm { C O _ { 2 } }$ 反应制备 $\mathrm { N a O C P }$ 合成⽅法。写出该制备过程的化学反应⽅程式。

\subsection*{2-4}
L $\mathrm { . i ( d m e ) _ { 2 } O C P }$ 与 $\mathrm { S O _ { 2 } }$ 反应时，OCP−会发⽣聚合，⽣成 $\mathrm { [ P _ { 4 } C _ { 4 } O _ { 4 } }$ ]2−，该离⼦包含两个五元环，且含P—P键， $\mathrm { ^ { 3 1 } P - N M R }$ 表明P原⼦有两种化学环境，画出 $\mathrm { [ P _ { 4 } C _ { 4 } O _ { 4 } ] }$ 2−的结构简式。

\subsection*{2-5}
NaOCP 与 $( \mathrm { i p c } ) _ { 2 } \mathrm { B C l }$ 反应，得到分⼦式为 $\mathrm { B _ { 3 } C _ { 3 } P _ { 3 } O _ { 3 } ( i p c ) _ { 6 } }$ 的化合物 $\mathbf { Y } _ { \circ }$ 。Y在叔丁醇的作⽤下，转变为化合物Z，Z有三重旋转轴，其31P-NMR测试表明所有P原⼦的化学环境相同。画出Y和Z的结构简式。

\subsection*{2-6}
研究发现： $\mathrm { O C P ^ { - } }$ 与过渡⾦属中⼼配位，会导致碳磷键断裂，⽣成CO分⼦。当⾦属中⼼有富余电⼦时，所⽣成的CO会与⾦属中⼼配位；⽽当⾦属中⼼缺电⼦时，CO会脱离反应体系。

${ \bf 2 - 6 - 1 } \left[ \mathrm { W ( O D i p p ) _ { 4 } } \right]$ 与NaOCP在⼄⼆醇⼆甲醚(dme)和12-冠-4存在下反应，得到产物 $\mathbf { Q } _ { \circ }$ 。写出 $[ \mathrm { W ( O D i p p ) _ { 4 } } ]$ 中W的氧化数，写出 $\mathrm { [ W ( O D i p p ) _ { 4 } }$ ]的⼏何构型，画出产物Q中的配阴离⼦的结构简式。

$2 { - } 6 { - } 2 \left[ \mathrm { I r C l ( P D I ^ { M e } ) } \right]$ 在低温下与NaOCP反应，也发⽣了碳磷键断裂，得到了双⾦属双磷化合物 $\mathbf { R } _ { \circ }$ 。写出$\mathrm { [ I r C l ( P D I ^ { M e } ) }$ ]中Ir的价电⼦组态和R的分⼦式（配体⽤缩写）及Ir的位数。

本题涉及的缩写：

（图：）
